% ===========================================================================
% 12-bit SAR ADC 桥接冗余CDAC双向递归位权校准系统研究报告
% 编译: xelatex -interaction=nonstopmode comprehensive_report.tex (×2)
% ===========================================================================
\documentclass[UTF8,a4paper,12pt]{ctexart}
\usepackage[margin=2.5cm]{geometry}
\usepackage{amsmath,amssymb,amsthm}
\usepackage{booktabs,longtable,array,tabularx}
\usepackage{graphicx}
\usepackage{listings,xcolor}
\usepackage{hyperref}
\usepackage{enumitem}
\usepackage{fancyhdr}
\usepackage{titlesec}

% --- hyperref ---
\hypersetup{colorlinks=true,linkcolor=blue!60!black,urlcolor=blue!60!black,
            citecolor=blue!60!black}

% --- 定理环境 ---
\theoremstyle{definition}
\newtheorem{theorem}{定理}[section]
\newtheorem{corollary}{推论}[theorem]
\newtheorem{lemma}{引理}[section]

% --- code style ---
\lstset{
  basicstyle=\ttfamily\footnotesize,
  breaklines=true,columns=fullflexible,
  frame=single,framerule=0.3pt,rulecolor=\color{gray!50},
  backgroundcolor=\color{gray!8},
  keywordstyle=\color{blue!60!black}\bfseries,
  commentstyle=\color{green!50!black},
  stringstyle=\color{red!60!black},
}

% --- 页眉页脚 ---
\pagestyle{fancy}
\fancyhf{}
\fancyhead[C]{\small 12-bit SAR ADC 桥接冗余CDAC双向递归位权校准系统}
\fancyfoot[C]{\thepage}
\renewcommand{\headrulewidth}{0.4pt}

% --- 章节格式 ---
\titleformat{\section}{\Large\bfseries\centering}{\thesection}{1em}{}
\titleformat{\subsection}{\large\bfseries}{\thesubsection}{1em}{}
\titleformat{\subsubsection}{\normalsize\bfseries}{\thesubsubsection}{1em}{}

% ===========================================================================
\begin{document}
% ===========================================================================

% --- 标题页 ---
\begin{titlepage}
\centering
\vspace*{3cm}
{\Huge\bfseries 12-bit SAR ADC 桥接冗余CDAC\\[6pt]双向递归位权校准系统研究报告}\\[1.5cm]
{\large 基于 Huang 2024 双向递归校准算法}\\[1cm]
{\large 版本：Huang V6（直接递归实现）}\\[0.5cm]
{\large 日期：2026年7月}\\[3cm]

\vfill
{\large SAR\_ADC\_IDEAL\_CAL\_REPLACEMENT\_V6\_CLEAN}\\[0.3cm]
{\large 项目文档编号：TR-2026-001}
\end{titlepage}

% --- 摘要 ---
\begin{abstract}
\noindent
本报告系统性地阐述一套面向 12-bit 桥接冗余 Split-CDAC SAR ADC 的双向递归位权
校准系统的设计与验证。项目历经 V5（架构设计阶段）、V6（3-target 原型实现）、
V6.2（发现 w2 re-correction 数学缺陷）、V6.3（缺陷修复与隐式 gauge 引入）、
V6.5（W5 粗调自动量程扩展）、V6.6（外部 anchor 通道合并）直至当前 Huang V6
（全直接递归二进制 CDAC 实现）的多次迭代演化。报告详细覆盖 CDAC 拓扑演化、
校准算法数学原理、版本历史与核心算法演化、数学可辨识性证明、方向与程度敏感性
分析、Dither 注入机制、文件清单与版本归档、仿真验证指南以及已知限制与未来工作。
所有结论均有严格的数学推导支撑（定理 1--6），并提供完整的 Verilog-A 与 RTL
实现参考。
\end{abstract}

\tableofcontents
\newpage

% ===================================================================
\section{项目背景与目录结构}
% ===================================================================

\subsection{项目概述}

本项目旨在为 12-bit 桥接冗余 Split-CDAC SAR ADC 设计并实现一套
Huang~2024\cite{huang2024} 风格的双向递归位权校准系统。该系统通过 D+/D$-$ 差分
测量消除比较器失调，利用低位 CDAC 电容作为校准 DAC（calDAC），从低位到高位
逐级递归测量并更新各电容的数字位权，从而在不依赖精密模拟匹配的前提下实现
高线性度。

项目当前阶段完成了从 V5 架构设计到 Huang V6 全直接递归实现的完整迭代，
形成了包含 Verilog-A 行为模型、可综合 RTL 代码、Spectre 仿真测试台、
严格数学证明文档在内的完整交付包。

\subsection{目录结构}

\begin{longtable}{p{3cm}p{10cm}}
\toprule
\textbf{目录} & \textbf{说明} \\
\midrule
\endhead
\bottomrule
\endfoot
\texttt{src/} & 核心 Verilog-A 模块：DEC\_CAL\_PHY（校准解码器）、SWITCH\_CAL（CDAC 开关阵列）、
SAR\_LOGIC\_0716（14-beat 非二进制 SAR 逻辑）、SYNC\_asnyc（时序生成器）、
COM\_ideal（理想比较器，含 CAL-gated dither） \\
\texttt{sim/} & 仿真辅助模型与顶层网表：tb\_huang\_calibration.scs（主测试台）、
cdac\_binary\_mc.va（二进制 CDAC 蒙特卡洛模型）、run\_mismatch\_sweep.sh（失配扫描脚本） \\
\texttt{rtl/} & 可综合 Verilog RTL 版本：DEC\_CAL\_PHY\_rtl.v（与 VA 引脚兼容、功能同步） \\
\texttt{doc/} & 文档与数学证明：HUANG\_RECURSIVE\_CALIBRATION\_DESIGN.md（完整设计文档）、
PROOF\_V6\_3\_IDENTIFIABILITY.md（严格数学证明）、STATIC\_CHECK\_REPORT.txt（静态检查报告） \\
\texttt{archive/} & 版本归档：V6、V6.2、V6.3、V6.5、V7 各版本的 VA 与 RTL 代码 \\
\texttt{ext/} & 扩展校准版本：DEC\_CAL\_PHY\_HUANG\_RECURSIVE\_V4.va（V4 完整递归参考实现） \\
\texttt{ref/} & 参考论文：Huang 2024 HKUST PhD Thesis \\
\bottomrule
\end{longtable}

\subsection{核心文件清单}

\begin{longtable}{p{4.5cm}p{3cm}p{7cm}}
\toprule
\textbf{文件} & \textbf{路径} & \textbf{说明} \\
\midrule
\endhead
\bottomrule
\endfoot
DEC\_CAL\_PHY.va & src/ & V6.6 校准解码器（3-target + W5 autoranging + external anchor） \\
DEC\_CAL\_PHY\_HUANG\_V5.va & src/ & V5 完整递归校准（9-target，H=122 设计文档） \\
DEC\_CAL\_PHY\_rtl.v & rtl/ & V6.6 可综合 RTL 镜像 \\
SWITCH\_CAL.va & src/ & CDAC 开关阵列（BITD\_CAL/BITU\_CAL 编码） \\
SAR\_LOGIC\_0716.va & src/ & 14-beat 非二进制 SAR 逻辑 \\
SYNC\_asnyc.va & src/ & 时序生成器 \\
COM\_ideal.va & src/ & 理想比较器（含 CAL-gated dither） \\
tb\_huang\_calibration.scs & sim/ & 参数化仿真测试台 \\
cdac\_binary\_mc.va & sim/ & 二进制 CDAC 蒙特卡洛模型 \\
\bottomrule
\end{longtable}

% ===================================================================
\section{CDAC 拓扑演化}
% ===================================================================

\subsection{原始设计：C$_B$=1C，H=128}

原始 CDAC 为全差分桥接结构：

\begin{itemize}
  \item \textbf{低位阵列}（通过桥接电容 C$_B$ 连接到顶板 P）：
    C$_{L1}$=1C, C$_{L2}$=2C, C$_{L3}$=4C, C$_{L4}$=6C, C$_{L5}$=10C,
    C$_{L6}$=16C, C$_{L7}$=24C，低位总电容 C$_{LT}$ = 63C
  \item \textbf{高位阵列}（直接连接到顶板 P）：
    C$_{H1}$=1C, C$_{H2}$=2C, C$_{H2R}$=2C（冗余）, C$_{H4}$=4C,
    C$_{H8}$=8C, C$_{H16}$=16C，高位总电容 C$_{HT}$ = 33C
  \item \textbf{桥接电容}：C$_B$ = 1C
  \item \textbf{N 侧}：镜像结构
\end{itemize}

\begin{theorem}[W$_{128}$ 不可校准]
在 C$_B$=1C（H=128）的 CDAC 上，无法使用 Huang 方法独立校准 W$_{128}$。
\end{theorem}

\textbf{证明}：桥接衰减比 $\alpha = (C_{LT} + C_B)/C_B = 64$，因此高位 1C 权重
$H = 2\alpha = 128$。calDAC（含 terminal）最大输出 = 低位总权重 + terminal = 
2$\times$63 + 1 = 127。校准 W$_{128}$ 需要残差 $R = W_{128} - W_{wall}$ 落在
calDAC 范围内。唯一的 wall 候选是低位全开（127），此时 $R = 1$。正方向注入时，
若 $W_{128}$ 有任何正失配，或比较器有正向失调，calDAC 输出 127 + 1 = 128
达到饱和，无法辨识。$\blacksquare$

更严重的是，由于桥接 DAC 的结构性质 $W_k = \Sigma(\text{lower}) + 1$（含 terminal），
所有级别的 gap 均为 0，整个递归链完全阻塞。

\subsection{候选 A：H=122}

为打破 gap=0 的限制，设计候选 A 将桥接电容调整为 C$_B$=1.05C：

\begin{equation}
  \alpha = \frac{63 + C_B}{C_B} = \frac{63 + 1.05}{1.05} = 61,\quad
  H = 2\alpha = 122
\end{equation}

\begin{longtable}{lcr}
\toprule
\textbf{物理电容} & \textbf{C 值} & \textbf{差分权重} \\
\midrule
\endhead
\bottomrule
\endfoot
高位 16C & 16 & 1952 \\
高位 8C & 8 & 976 \\
高位 4C & 4 & 488 \\
高位 2C\_A & 2 & 244 \\
高位 2C\_R & 2 & 244 \\
高位 1C & 1 & 122 \\
低位 24C & 24 & 48 \\
低位 16C & 16 & 32 \\
低位 10C & 10 & 20 \\
低位 6C & 6 & 12 \\
低位 4C & 4 & 8 \\
低位 2C & 2 & 4 \\
低位 1C & 1 & 2 \\
terminal & --- & 1 \\
\bottomrule
\end{longtable}

总权重 = 4153，冗余 = 58。每级获得 5 LSB 的校准裕量，支持 9-target 完整递归。

\subsection{最终二进制 CDAC：C$_B$=2C，H=65}

当前 Huang V6 实际使用的系统采用二进制 CDAC 拓扑：

\begin{itemize}
  \item \textbf{低位阵列}：\{1, 2, 4, 8, 16, 32\} = 63C
  \item \textbf{桥接电容}：C$_B$ = 2C
  \item \textbf{高位阵列}：\{1A, 1R, 2, 4, 8, 16, 32\} = 64C
\end{itemize}

桥接衰减比 $\alpha = (63+2)/2 = 32.5$，高位单位权重 H = 65。

\textbf{权重表与映射}：

\begin{longtable}{ccl}
\toprule
\textbf{数字权重} & \textbf{映射关系} & \textbf{说明} \\
\midrule
\endhead
\bottomrule
\endfoot
2080 & W$_0$ = 32C $\times$ 65 & \\
1040 & W$_1$ = 16C $\times$ 65 & \\
520 & W$_2$ = 8C $\times$ 65 & \\
260 & W$_3$ = 4C $\times$ 65 & \\
130 & W$_4$ = 2C $\times$ 65 & \\
65  & high 1C\_A & stage 5 \\
65  & high 1C\_R & stage 6（冗余）\\
64  & low 32C $\times$ 2 & \\
32  & low 16C $\times$ 2 & \\
16  & low 8C $\times$ 2 & \\
8   & low 4C $\times$ 2 & \\
4   & low 2C $\times$ 2 & \\
2   & low 1C $\times$ 2 & \\
1   & terminal & \\
\bottomrule
\end{longtable}

总权重 = 2080+1040+520+260+130+65+65+64+32+16+8+4+2+1 = \textbf{4287}。
冗余 = 4287 $-$ 4095 = 192。

\textbf{Stage 映射}：

\begin{longtable}{clc}
\toprule
\textbf{Stage} & \textbf{电容} & \textbf{权重} \\
\midrule
\endhead
\bottomrule
\endfoot
0 & H32C & 2080 \\
1 & H16C & 1040 \\
2 & H8C & 520 \\
3 & H4C & 260 \\
4 & H2C\_A & 130 \\
5 & H1C\_A & 65 \\
6 & H1C\_R & 65 \\
7 & L32C & 64 \\
8 & L16C & 32 \\
9 & L8C & 16 \\
10 & L4C & 8 \\
11 & L2C & 4 \\
12 & L1C & 2 \\
13 & terminal & 1 \\
\bottomrule
\end{longtable}

% ===================================================================
\section{版本历史与核心算法演化}
% ===================================================================

\subsection{版本总览}

\begin{longtable}{p{2cm}p{3cm}p{9cm}}
\toprule
\textbf{版本} & \textbf{状态} & \textbf{核心特征} \\
\midrule
\endhead
\bottomrule
\endfoot
V5 & 设计文档 & 9-target 完整递归，H=122 CDAC，完整数学推导与切换表。未实现代码。 \\
V6 & 归档 & 3-target 原型，7-bit calDAC（$\pm 126$），E=e3+e4 传播。无饱和检测。 \\
V6.1 & (未归档) & 过渡版本，引入基本饱和检测。 \\
V6.2 & DEFECT & w2 re-correction 引入 bug：雅可比矩阵奇异，修正为恒等操作。 \\
V6.3 & 归档 & 删除 w2 re-correction，引入隐式 gauge（W3+W4=512），修复 P0/P1 bugs。 \\
V6.5 & 草案 & W5 粗调自动量程扩展（$\pm 122 \to \pm 250$）。ChatGPT 生成。 \\
V6.6 & 当前发布 & 合并 V6.5 自动量程 + 外部 anchor 通道（§7 方案 A）。 \\
V7 & DEFECTIVE & 128>126 违反 offset-binary SAR 收敛。请勿使用。 \\
Huang V6 & 当前实现 & 全直接递归，二进制 CDAC 适配，7-target，可变 search\_count（6..12）。 \\
\bottomrule
\end{longtable}

\subsection{V5：9-target 完整递归（架构设计）}

V5 是项目的架构设计阶段，基于候选 A（H=122）CDAC，规划了从 low10（20）到
high16（1952）共 9 个目标的完整递归校准序列。每级 calDAC 配置、切换表、
误差传播模型均在设计文档（doc/HUANG\_RECURSIVE\_CALIBRATION\_DESIGN.md）
中完整推导。V5 未实现可运行代码，但为后续所有版本提供了理论基础。

\subsection{V6：3-target 原型}

V6 实现 3-target 校准（W$_2$=512, W$_1$=1024, W$_0$=2048），使用 7-bit calDAC
（$\pm 126$）。核心特征：
\begin{itemize}
  \item 3-target 序列：target 2 (W$_2$) $\to$ target 1 (W$_1$) $\to$ target 0 (W$_0$)
  \item 隐式 gauge：W$_3$+W$_4$=512（名义值），不单独校准
  \item E = e$_3$+e$_4$ 沿 $(-E, -2E, -4E)$ 传播到 (w$_2$, w$_1$, w$_0$)
  \item 无饱和检测，无 P0/P1 修复
\end{itemize}

\subsection{V6.2：w2 re-correction 缺陷}

V6.2 在 V6 基础上试图通过"w2 re-correction"修正 E 传播：
\begin{lstlisting}[language=verilog]
w2 = w2 + (measured_weight_q - 512*q_scale);
\end{lstlisting}

\begin{theorem}[V6.2 w2 re-correction 是 no-op]
在 V6.2 的 4-target 序列下，第 2 次测量必然得到 $\hat W_{34} = 512$（与 e$_2$, E 无关），
因此 w2 re-correction 是恒等操作。
\end{theorem}

\textbf{证明}：第一次测量 $\hat W_2 = 512 + (e_2 - E)$。第二次测量 target = W$_{34}$，
wall = $\hat W_2$，R$_2$ = E $-$ e$_2$。测量结果 $\hat W_{34} = \hat W_2 + R_2 = 512$。
雅可比矩阵 $\det(J) = 0$，参数 (e$_2$, E) 不可辨识。$\blacksquare$

\subsection{V6.3：缺陷修复与隐式 gauge}

V6.3 删除 V6.2 的 no-op 操作，修复 P0/P1 系列 bug：
\begin{itemize}
  \item \textbf{P0-1}：wall\_weight\_q 初始化（三处设置 wall = w3+w4）
  \item \textbf{P0-2}：RST1 双 handler 合并为单一 handler
  \item \textbf{P0-3}：REJECT 静默继续 $\to$ fail-fast ERR=1
  \item \textbf{P1-1}：saturate\_flag 单 bit $\to$ per-direction 计数
  \item \textbf{P1-2}：saturated sample 排除 accum
  \item \textbf{P1-3}：死参数清理
  \item \textbf{P1-4}：decoder 在 ERR=1 时回退名义权重
\end{itemize}

\subsection{V6.5：W5 粗调自动量程}

V6.5 解决 R 量程饱和问题（根因 II）。将 W$_5$=128 作为粗调电容，通过自动量程
选择将 calDAC 可测范围从 $\pm 122$ 扩展至 $\pm 250$ LSB。

\begin{theorem}[W5 自动量程扩展]
在 W5 粗调自动量程下，可测对称残差范围从 $\pm 122$ LSB 扩展到 $\pm 250$ LSB。
\end{theorem}

\textbf{证明}：calDAC 安全测量范围 $|R_{fine}| \leq 122$。k=$\pm 1$ 时加入
$\pm 128$ 粗调偏置，覆盖 $R \in [-250, -6] \cup [-122, 122] \cup [6, 250]$。$\blacksquare$

\subsection{V6.6：最终合并版本}

V6.6 合并 V6.5 自动量程与外部 anchor 通道（§7 方案 A），同时解决两个独立根因。

\textbf{外部 Anchor 通道}：暴露 ANCHOR\_W3\_LSB 和 ANCHOR\_W4\_LSB 参数。

\begin{theorem}[V6.6 anchor 误差传播]
在 V6.6 3-target 递归下：
\begin{align*}
\hat W_2 &= 512 + e_2 - E + e_{REF} \\
\hat W_1 &= 1024 + e_1 - 2E + 2e_{REF} \\
\hat W_0 &= 2048 + e_0 - 4E + 4e_{REF}
\end{align*}
即 anchor 误差 e$_{REF}$ 按 $(+e_{REF}, +2e_{REF}, +4e_{REF})$ 传播，符号与 E 相反。
\end{theorem}

当用户提供精确物理表征值使 e$_{REF}$ = E 时，E 传播完全消除。

\subsection{Huang V6（当前）：二进制 CDAC 全直接递归}

当前实现版本采用二进制 CDAC 拓扑（C$_B$=2C，H=65），实现全直接递归校准：
\begin{itemize}
  \item 7-target 校准序列
  \item 可变 search\_count：各目标 6..12 步
  \item D+/D$-$ 32-pair 平均
  \item Q6 固定小数位格式
\end{itemize}

% ===================================================================
\section{校准算法原理}
% ===================================================================

\subsection{Huang 双向递归方法}

Huang 2024 提出的双向递归位权校准方法的核心思想是：利用低位已校准电容作为
calDAC，通过 D+/D$-$ 双向差分测量消除比较器失调，逐级递归测量高位电容的
实际物理权重。

\textbf{校准序列}（以二进制 CDAC 为例）：
\begin{align*}
\text{target 6: } &\hat W_7 = 64 + R_7 \quad (\text{wall} = W_{8..13} + W_6) \\
\text{target 5: } &\hat W_6 = \hat W_7 + \text{wall}_{6} + R_6 \\
&\vdots \\
\text{target 0: } &\hat W_0 = \hat W_1 + \hat W_2 + \cdots + \text{wall}_0 + R_0
\end{align*}

\subsection{D+/D$-$ 方法消除 V$_{OS}$}

对每个待校准位 W$_k$，执行双向测量：
\begin{align*}
D^+ &= W_k + V_{OS} + x_n \quad (\text{正方向注入}) \\
D^- &= -W_k + V_{OS} + x_{n+1} \quad (\text{负方向注入})
\end{align*}

位权估计：
\begin{equation}
\hat W_k = \frac{D^+ - D^-}{2} = W_k + \frac{x_n - x_{n+1}}{2}
\end{equation}

V$_{OS}$ 完全消除，仅残留比较器噪声差。

\subsection{Wall-residual 与直接递归对比}

\textbf{Wall-residual 方法}（V6.3/V6.6）：使用 W$_3$+W$_4$ 作为 wall 的
名义值（512），测量残差 R = W$_{target}$ $-$ W$_{wall}$。数字侧记录
$\hat W_{target}$ = wall$_{digital}$ + R。

\textbf{直接递归方法}（Huang V6/V5）：每个 target 直接使用之前已校准的
低位权重作为 calDAC，不依赖固定的 wall 名义值。calDAC 范围和 SAR 搜索
步数逐级增加。

\subsection{32-pair 平均}

校准重复 32 个 D+/D$-$ pair：
\begin{equation}
\hat W_k = \frac{1}{32} \sum_{m=1}^{32} \frac{D^{+(m)} + |D^{-(m)}|}{2}
= \frac{\text{accum\_plus} + \text{accum\_minus\_mag}}{2 \times 32}
\end{equation}

噪声标准差降低 8 倍：
\begin{equation}
\sigma_{\hat W} = \frac{\sigma_n}{\sqrt{32}} = \frac{\sigma_n}{8}
\end{equation}

\subsection{当前二进制 CDAC 的 search\_count}

各目标 SAR 搜索步数（含 terminal）：

\begin{longtable}{ccc}
\toprule
\textbf{目标} & \textbf{Stage} & \textbf{search\_count} \\
\midrule
\endhead
\bottomrule
\endfoot
W$_{64}$ (calDAC) & 7 & 6 \\
W$_{65}$A & 5 & 7 \\
W$_{65}$R & 6 & 7 \\
W$_{130}$ & 4 & 8 \\
W$_{260}$ & 3 & 9 \\
W$_{520}$ & 2 & 10 \\
W$_{1040}$ & 1 & 11 \\
W$_{2080}$ & 0 & 12 \\
\bottomrule
\end{longtable}

\subsection{Trial stage 映射公式}

对于 target 位于 stage $s_t$，calDAC 包含 stages $\{s_{t+1}, s_{t+2}, \ldots, s_{13}\}$
以及已校准的高位 stages（在校准序列中位于 target 之后的目标）。SAR trial 顺序
从大到小。

% ===================================================================
\section{数学可辨识性证明（摘要）}
% ===================================================================

本节摘录项目中六个核心定理的证明要点。完整推导见
doc/PROOF\_V6\_3\_IDENTIFIABILITY.md。

\begin{theorem}[V6.2 w2 re-correction 是 no-op]
在 V6.2 的 4-target 序列下，$\hat W_{34} \equiv 512$，与 e$_2$, E 无关。
雅可比矩阵 $\det(J) = 0$，参数 (e$_2$, E) 不可辨识。
\end{theorem}

\begin{theorem}[W5 (128) 不可测]
用现有 calDAC ($\pm 126$) 无法测量 W$_5$ = 128 + e$_5$。
calDAC 与 wall 共享物理电容，唯一的 wall 候选是 W$_{13}$=1，
R = 127 > 126。
\end{theorem}

\begin{theorem}[共模增益不可观测]
所有 wall 测量对共模增益 (1+$\alpha$) 不敏感。calDAC 自身也缩放，
数字 code 与 $\alpha$ 无关。
\end{theorem}

\begin{theorem}[V6.3 E 传播（更正版）]
在 V6.3 3-target 序列下：
\begin{align*}
\hat W_2 &= 512 + e_2 - E \\
\hat W_1 &= 1024 + e_1 - 2E \\
\hat W_0 &= 2048 + e_0 - 4E
\end{align*}
E 按 $(-E, -2E, -4E)$ 传播，不是精确校准。
\end{theorem}

\begin{theorem}[W5 自动量程扩展]
在 W5 粗调自动量程下，可测对称残差范围从 $\pm 122$ 扩展到 $\pm 250$ LSB。
覆盖 $\pm 5\%$ 任意方向 MSB 失配，$\pm 7\%$ 边界，$\pm 10\%$ 仍失败。
\end{theorem}

\begin{theorem}[Anchor e$_{REF}$ 传播]
在 V6.6 3-target 递归下，anchor 误差 e$_{REF}$ 按
$(+e_{REF}, +2e_{REF}, +4e_{REF})$ 传播。当 e$_{REF}$ = E 时，E 完全消除。
\end{theorem}

% ===================================================================
\section{方向与程度敏感性分析}
% ===================================================================

\subsection{两个独立根因}

V6.3 在不同失配方向上校准效果不一致有两个独立根因：

\begin{longtable}{p{2.5cm}p{4cm}p{3.5cm}p{4cm}}
\toprule
\textbf{根因} & \textbf{描述} & \textbf{触发条件} & \textbf{V6.6 解决方案} \\
\midrule
\endhead
\bottomrule
\endfoot
(I) E 传播 & E=e$_3$+e$_4$ 作为 $(-E,-2E,-4E)$ 进入 (w$_2$,w$_1$,w$_0$) & 任何 E $\neq$ 0 & 外部 anchor A$_{34}$ = W$_3$+W$_4$（物理表征）\\
(II) R 量程饱和 & R$_0$=e$_0$-e$_1$-e$_2$ 超过 $\pm 126$ LSB & 反向失配且 |p| $>$ 3.5\% & W5 粗调自动量程，扩展至 $\pm 250$ LSB \\
\bottomrule
\end{longtable}

\subsection{方向敏感性}

设 |e$_3$| = |e$_4$| = $\sigma$（相同程度），方向不同：

\begin{longtable}{lccccc}
\toprule
\textbf{方向} & \textbf{e$_3$} & \textbf{e$_4$} & \textbf{E} & $\Delta$W$_0$ = $-4$E & \textbf{效果} \\
\midrule
\endhead
\bottomrule
\endfoot
反向（对消） & $+\sigma$ & $-\sigma$ & 0 & 0 & $\checkmark$ 完美 \\
同向（累加） & $+\sigma$ & $+\sigma$ & $+2\sigma$ & $-8\sigma$ & $\times$ 失败 \\
同向（反向） & $-\sigma$ & $-\sigma$ & $-2\sigma$ & $+8\sigma$ & $\times$ 失败 \\
\bottomrule
\end{longtable}

\subsection{程度敏感性}

\begin{longtable}{ccccccc}
\toprule
\textbf{|E|} & \textbf{$|\Delta$W$_0|$} & \textbf{$|\Delta$W$_1|$} & \textbf{$|\Delta$W$_2|$} & \textbf{最坏 INL} & \textbf{ENOB} \\
\midrule
\endhead
\bottomrule
\endfoot
0.512 & 2.05 & 1.02 & 0.51 & $\sim 3.6$ & $\sim 11.5$ \\
2.56 & 10.24 & 5.12 & 2.56 & $\sim 18$ & $\sim 10.5$ \\
5.12 & 20.48 & 10.24 & 5.12 & $\sim 36$ & $\sim 9.5$ \\
25.6 & 102.4 & 51.2 & 25.6 & $\sim 179$ & $\sim 7.5$ \\
\bottomrule
\end{longtable}

\subsection{外部参考电容方案（Solution A）}

引入独立参考电容 C$_{REF}$=512（与 W$_3$+W$_4$ 名义相同），物理隔离。
校准序列：

\begin{align*}
\hat W_2 &= 512 + e_2 - e_{REF} \\
\hat W_1 &= 1024 + e_1 - 2e_{REF} \\
\hat W_0 &= 2048 + e_0 - 4e_{REF}
\end{align*}

E 被 e$_{REF}$ 替换。e$_{REF}$ 是单个独立电容的失配，不存在"同向累加"问题，
彻底消除方向敏感性。

% ===================================================================
\section{DITHER 注入分析}
% ===================================================================

\subsection{Dither 设计参数}

\begin{longtable}{lcl}
\toprule
\textbf{参数} & \textbf{值} & \textbf{说明} \\
\midrule
\endhead
\bottomrule
\endfoot
Peak A & $\pm 0.5$ calibration LSB & 亚 LSB 可观测性 \\
序列 & 32-level deterministic uniform & 确定性均匀序列 \\
单位 & terminal LSB (= 1) & 非物理 1C weight=2 \\
Q6 匹配 & 1/64 LSB resolution & 6-bit 小数精度 \\
更新方式 & Per-pair update & 同一 dither 用于 D+ 和 D$-$ \\
\bottomrule
\end{longtable}

\subsection{工作原理}

Dither 信号由 COM\_ideal 模块中的 CAL 引脚控制，仅在校准模式开启。
对于晶体管级比较器（含实际噪声），DITHER 不需要——实际比较器热噪声
已提供天然的亚 LSB 可观测性。DITHER 是仿真辅助机制。

Per-pair 更新保证 D+ 和 D$-$ 使用相同的 dither 值，因此差分抵消后
dither 贡献被消除。

% ===================================================================
\section{文件清单与版本归档}
% ===================================================================

\subsection{src/ 文件清单}

\begin{longtable}{p{5cm}p{9cm}}
\toprule
\textbf{文件} & \textbf{说明} \\
\midrule
\endhead
\bottomrule
\endfoot
DEC\_CAL\_PHY.va & V6.6 校准解码器（3-target + W5 autoranging + external anchor） \\
DEC\_CAL\_PHY\_HUANG\_V5.va & V5 完整递归校准（9-target，H=122 架构设计） \\
DEC\_CAL\_PHY\_HUANG\_V6.va & Huang V6 直接递归实现（二进制 CDAC，7-target） \\
SWITCH\_CAL.va & CDAC 开关阵列（BITD\_CAL/BITU\_CAL 编码） \\
SAR\_LOGIC\_0716.va & 14-beat 非二进制 SAR 逻辑 \\
SYNC\_asnyc.va & 时序生成器 \\
COM\_ideal.va & 理想比较器（含 CAL-gated dither） \\
\bottomrule
\end{longtable}

\subsection{sim/ 文件清单}

\begin{longtable}{p{5cm}p{9cm}}
\toprule
\textbf{文件} & \textbf{说明} \\
\midrule
\endhead
\bottomrule
\endfoot
tb\_huang\_calibration.scs & 主测试台（参数化，支持多种失配注入模式） \\
cdac\_binary\_mc.va & 二进制 CDAC 蒙特卡洛模型 \\
cdac\_split\_mc\_thesis\_deltaC.va & 论文版 split-CDAC MC 模型 \\
cdac\_split\_mc\_final\_1pct.va & 最终 split-CDAC MC 模型（1\% $\sigma$） \\
run\_mismatch\_sweep.sh & 失配扫描自动化脚本 \\
TOP\_LEVEL\_CONNECTION\_PATCH.scs & 顶层连接补丁 \\
SIMULATION\_GUIDE.md & 仿真指南 \\
\bottomrule
\end{longtable}

\subsection{archive/ 版本归档}

\begin{longtable}{p{3cm}p{4cm}p{7cm}}
\toprule
\textbf{版本} & \textbf{文件} & \textbf{说明} \\
\midrule
\endhead
\bottomrule
\endfoot
V6 & DEC\_CAL\_PHY\_v6.va & 3-target，7-bit calDAC，无饱和检测 \\
V6.2 & DEC\_CAL\_PHY\_v6\_2.va & DEFECT：w2 re-correction 数学无效 \\
V6.2 RTL & DEC\_CAL\_PHY\_rtl\_v6\_2.v & V6.2 RTL 镜像 \\
V6.3 & DEC\_CAL\_PHY\_v6\_3.va & 3-target + 隐式 gauge + P0/P1 修复 \\
V6.3 RTL & DEC\_CAL\_PHY\_rtl\_v6\_3.v & V6.3 RTL 镜像 \\
V6.5 & DEC\_CAL\_PHY\_v6\_5.va & W5 粗调自动量程草案 \\
V7 & DEC\_CAL\_PHY\_v7.va & DEFECTIVE：128$>$126 违反 offset-binary SAR \\
\bottomrule
\end{longtable}

\subsection{doc/ 文件清单}

\begin{longtable}{p{5cm}p{9cm}}
\toprule
\textbf{文件} & \textbf{说明} \\
\midrule
\endhead
\bottomrule
\endfoot
README\_使用说明.md & 项目使用说明（V6.6 完整文档） \\
PROOF\_V6\_3\_IDENTIFIABILITY.md & 严格数学证明（定理 1--6 + §7 文献综述 + §8 V6.6 证明） \\
HUANG\_RECURSIVE\_CALIBRATION\_DESIGN.md & V5 完整设计文档（电荷方程、拓扑演化、误差传播、伪代码） \\
STATIC\_CHECK\_REPORT.txt & V6.6 静态检查报告 \\
sar\_adc\_mismatch\_calibration\_report\_cn.pdf & V6.2 时期报告（历史参考） \\
comprehensive\_report.tex & 本报告 \\
\bottomrule
\end{longtable}

% ===================================================================
\section{仿真验证指南}
% ===================================================================

\subsection{Testbench 结构}

主测试台为 \texttt{sim/tb\_huang\_calibration.scs}，包含：
\begin{itemize}
  \item CDAC 模型（参数化失配注入）
  \item SWITCH\_CAL 开关阵列（无需修改）
  \item SAR\_LOGIC\_0716 SAR 逻辑（无需修改）
  \item DEC\_CAL\_PHY 校准解码器（当前版本）
  \item COM\_ideal 理想比较器（含 CAL-gated dither）
  \item SYNC\_asnyc 时序生成器
  \item 激励源：CAL\_RST, RST1, CAL\_CLK, input ramp/sine
\end{itemize}

\subsection{测试用例}

\begin{longtable}{p{2.5cm}p{4cm}p{3cm}p{4.5cm}}
\toprule
\textbf{类别} & \textbf{名称} & \textbf{CDAC} & \textbf{验证项} \\
\midrule
\endhead
\bottomrule
\endfoot
Ideal & ideal\_all & ideal & $\hat W_i = W_i$，误差仅来自有限平均 \\
Single mismatch & single\_cap & 单电容 $\pm 5\%$ & $\hat W_i$ 失配检测精度 \\
Dual-side & pn\_same/oppo & P/N 同/反向 & 差分权重正确性 \\
MC & mc\_3msb/mc\_all & 全局 MC (1\% $\sigma$) & INL/SNDR/ENOB 统计分布 \\
\bottomrule
\end{longtable}

\subsection{关键参数}

\begin{longtable}{lcl}
\toprule
\textbf{参数} & \textbf{默认值} & \textbf{说明} \\
\midrule
\endhead
\bottomrule
\endfoot
NOM\_H1\_WEIGHT & 65 & 高位单位权重（二进制 CDAC） \\
AVG\_PAIRS\_LOG2 & 5 & 32 对 D+/D$-$ 平均 \\
FRAC\_BITS & 6 & Q6 固定小数位 \\
CAL\_BYPASS & 0 & 0=运行校准，1=旁路 \\
APPLY\_CAL\_WEIGHT & 1 & 是否应用校准权重 \\
WEIGHT\_TOL\_PCT & 35 & 位权 35\% 容差窗 \\
\bottomrule
\end{longtable}

\subsection{预期输出格式}

\begin{lstlisting}
HUANGV6 init ... targets=7
HUANGV6 measure target=6 ... ACCEPT valid=32/32 sat=0/32 (w7)
HUANGV6 measure target=5 ... ACCEPT valid=32/32 sat=0/32 (w6)
...
HUANGV6 finish done=1 err=0
HUANGV6 weights_q={2080,1040,520,260,130,65,65,64,32,16,8,4,2,1}
\end{lstlisting}

% ===================================================================
\section{已知限制与未来工作}
% ===================================================================

\subsection{已知限制}

\begin{enumerate}
  \item \textbf{W$_5$ 不可校准}：定理 2 证明 128 不可测（唯一 wall 候选 W$_{13}$=1，
    R=127$>$126）。当前依赖 layout matching 保证 $|e_5| < 1\%$。
  \item \textbf{e$_3$/e$_4$ 不可分离}：只假设 W$_3$+W$_4$=512（gauge）。
    e$_3$$-$e$_4$ 差分失配作为 DNL 噪声保留。
  \item \textbf{$\pm 7\%$ 失配边界}：5\% 任意方向覆盖良好但 7\% 接近量程边界
    （$|R_0|_{max}=250.9$）。需第 2 级粗调电容覆盖 $\pm 10\%$。
  \item \textbf{未进行 Spectre/AMS 仿真}：当前仅为 Verilog-A 行为级验证，
    尚未在 Spectre 晶体管级平台上运行完整仿真。
  \item \textbf{共模增益不可观测}：定理 3 证明所有 wall 测量对共模增益不敏感。
    整体增益误差需要外部 V$_{ref}$ 校准。
  \item \textbf{方向歧义 fail-fast}：若 k=0 探测同时触发 overflow\_pos 和
    overflow\_neg，立即设 ERR=1。
\end{enumerate}

\subsection{未来工作}

\begin{enumerate}
  \item \textbf{Spectre 仿真验证}：在晶体管级 CDAC 和比较器模型上运行完整
    校准仿真，验证算法在真实电路噪声环境下的收敛性。
  \item \textbf{Monte Carlo 验证}：100+ 次 MC 仿真，统计 ENOB 分布与
    最坏情况 INL。
  \item \textbf{W$_5$ 硬件修复}：将 W$_5$=128 拆分为 64+64，校准子电容；
    或引入外部参考电容作为 wall。
  \item \textbf{二级粗调电容}：利用 W$_4$=256 作为二级粗调，将量程扩展至
    $\pm 506$ LSB，覆盖 $\pm 10\%$ 失配。
  \item \textbf{外部 Anchor 硬件实现}：集成 C$_{REF}$ 电容与切换逻辑，
    实现完整的物理 anchor 通道。
  \item \textbf{Split-ADC 后台校准}：McNeill 2005 架构，两个半尺寸 ADC
    互为参考，彻底消除递归自引用。
  \item \textbf{两段式校准}：Chen 2024 TCAS-I 启发，将 MSB 段与 LSB 段
    独立校准，避免 E 传播。
\end{enumerate}

% ===================================================================
\section*{参考文献}
% ===================================================================
\addcontentsline{toc}{section}{参考文献}

\begin{thebibliography}{99}

\bibitem{huang2024}
Q. Huang,
\textit{Advanced Clock Multiplier and SAR ADC Design Techniques for
High-Resolution Signal Chain Systems},
HKUST PhD Thesis, 2024, Section 4.5.

\bibitem{lee1984}
H.-S. Lee, D. A. Hodges, P. R. Gray,
``A Self-Calibrating 15 bit CMOS A/D Converter,''
\textit{IEEE J. Solid-State Circuits}, vol.~SC-19, no.~6, pp.~813--819, 1984.

\bibitem{chen2024tcas1}
Y. Chen, Q. Huang, Y. Fan, Q. Zhao, S. Huang, J. Yuan,
``A 16-bit 4-MS/s SAR ADC With Dual-Segmental Bit Weight Self-Calibration,''
\textit{IEEE Trans. Circuits Syst. I}, vol.~71, no.~9, pp.~3961--3974, 2024.

\bibitem{chen2024tvlsi}
Y. Chen, S. Huang, Q. Huang, Y. Fan, J. Yuan,
``The Error Analysis of Bit Weight Self-Calibration Methods for
High-Resolution SAR ADCs,''
\textit{IEEE Trans. VLSI Syst.}, vol.~32, no.~11, pp.~1983--1992, 2024.

\bibitem{fan2025}
H. Fan, Z. Chen, Y. Liu, F. Maloberti, M. Chen, Q. Wei,
``Analog Foreground Calibration of High-Resolution SAR ADC,''
\textit{IEEE Trans. Circuits Syst. I}, vol.~72, no.~11, pp.~6629--6639, 2025.

\bibitem{mcneill2005}
J. A. McNeill, M. C. W. Coln, B. J. Larivee,
``'Split ADC' Architecture for Deterministic Digital Background Calibration
of a 16-bit 1-MS/s ADC,''
\textit{IEEE J. Solid-State Circuits}, vol.~40, no.~12, pp.~2437--2445, 2005.

\bibitem{um2013}
J.-Y. Um, Y.-J. Kim, E.-W. Song, J.-Y. Sim, H.-J. Park,
``A Digital-Domain Calibration of Split-Capacitor DAC for a Differential SAR ADC
Without Additional Analog Circuits,''
\textit{IEEE Trans. Circuits Syst. I}, vol.~60, no.~11, pp.~2845--2856, 2013.

\bibitem{ginsburg2005}
B. P. Ginsburg, A. P. Chandrakasan,
``An Energy-Efficient Charge Recycling Approach for a SAR Converter
with Capacitive DAC,''
\textit{IEEE ISCAS 2005}, pp.~184--187.

\bibitem{lopez2020}
A. L\'{o}pez-Angulo, A. J. Gin\'{e}s, E. J. Peralias,
``Digital calibration of capacitor mismatch and comparison offset in
Split-CDAC SAR ADCs with redundancy,''
\textit{IEEE ISCAS 2020}.

\bibitem{ding2025}
Y. Ding et al.,
``A 16-bit 1-MS/s SAR ADC With Capacitor Mismatch Self-Calibration,''
2025.

\bibitem{zheng2024}
Y. Zheng,
``小尺寸电容阵列高精度逐次逼近型ADC研究,''
2024.

\end{thebibliography}

% ===========================================================================
\end{document}
